Explanation methods such as Grad-CAM, RISE and LIME are increasingly attached
to deepfake detectors so that a forensic analyst can see \emph{why} an image was
flagged as synthetic. Whether those explanations can be trusted is usually
judged with perturbation-based faithfulness metrics: the pixels an explanation
marks as important are deleted, and a faithful explanation should make the
detector's score fall quickly. The metric, however, does not say what to put
in place of the deleted pixels. That choice is left to a \emph{fill operator}
(black, grey, blur, noise and so on). This project starts from the observation
that, for a forgery detector, the fill operator is not a harmless detail. Hard
mask edges, flat grey patches and white noise are exactly the kind of
low-level artifacts a forgery detector is trained to notice, so the measuring
instrument may itself produce the signal it is supposed to measure.

To test this, we designed a \emph{leakage null test} that needs no explanation
method at all: random regions of genuinely real photographs are masked and
filled, and any rise in the detector's probability of ``fake'' is attributed to
the operator. We ran this test, together with a leakage-corrected faithfulness
score ($\Delta$AUC, salient deletion measured against a random-deletion
control), on three frozen backbones (ResNet-50, ViT-B/16 and CLIP ViT-B/16)
with logistic-regression heads, about 171\,000 images from roughly forty
generators grouped into six families, eight fill operators and seven
explanation methods.

The results show that the instrument leaks heavily. Random masking of real
images raises CLIP's mean probability of ``fake'' by up to 0.62 under grey
fill, and every detector has at least one standard operator with leakage above
0.2. No operator is neutral for all detectors: transplanting pixels from
another real photograph is nearly leak-free on ResNet-50 (0.12) and ViT-B/16
(0.002), but still leaks 0.58 on CLIP. The fill operator also reorders
explanation methods: on ViT-B/16 the rank correlation between two operators
drops to $-0.29$, and gradient$\times$input moves from first to last place.
LIME stays at or near the top in every setting, so the main conclusion of
earlier work survives, but full method rankings measured with a single
operator do not. Finally, in a controlled experiment where only the generator
diversity of the training set is changed, diversity consistently improves
detection AUC on unseen families (+0.04 to +0.05), while the change in
faithfulness depends strongly on the operator, and faithfulness does not track
detection accuracy in general. These findings argue that faithfulness numbers
for forensic detectors should always be reported together with the leakage of
the operator used to compute them.

\vspace{0.5cm}

